\subsection{Finite-order estimates for the exact long passage}
\label{rw:app:secondary-long}

The exact saddle power accounts for the nonintegrable part of the long
flight. Its remaining correction is integrable, and differentiating
that correction proves Proposition~\ref{rw:exact-long}, including the
ordinary terminal-height jets at zero. Fix the total order $N$ first.
The constants may depend on $N$ and the fixed heights and sections,
but are independent of the long time.

Use the terminal variables of \eqref{rw:secondary:long-scalar}.
The rectangle is chosen with $W<-c$ and $Q>1/2$, so
$0\le w(t)\le w_e e^{-t/2}$. The exact subtraction of the saddle
ratio gives
\begin{equation}\label{rw:secondary:integrable-carriers}
 |Q-\beta_\epsilon|\le C\left(
 |\gamma|e^{-t}+|J|e^{-2t}+w+
 \epsilon(|\bar a|+|\nu|)\varphi w+\epsilon\varphi^2w\right).
\end{equation}
A coefficient derivative may remove the small coefficient multiplying
a term, but leaves its exponential or weighted factor. These factors
control the differentiated integral: for every fixed integer $i$,
\[
 \int_0^L(1+t)^iw(t)\,\dd t\le C_i,\qquad
 \int_0^L(1+t)^i(\varphi+\varphi^2)w(t)\,\dd t
       \le C_iH^{1/2}(1+L)^i .
\]
Indeed $\varphi=He^t$ and $e^L=\varphi_0/H$; integration of
$He^{t/2}$ and $H^2e^{3t/2}$ proves the second estimate.

Write
\[
 w=w_e e^{-\beta_\epsilon t}e^{R(t)},\qquad R(0)=0,\qquad
 R_t=-(Q-\beta_\epsilon).
\]
The zeroth-order bound follows by integrating
\eqref{rw:secondary:integrable-carriers}.
At each derivative order, the highest variation of $R$ solves a linear
inhomogeneous equation with coefficient $-wQ_w$. Its absolute integral
is bounded independently of $L$, so variation of constants has a common
bound. A $\gamma$ or $J$ derivative introduces $e^{-t}$ or $e^{-2t}$;
an $\bar a$ or $\nu$ derivative introduces $\epsilon\varphi w$.
An $H\partial_H$ derivative of an explicit coefficient retains
$\varphi w$ or $\varphi^2w$. A $\theta$ derivative of the removed
power introduces a factor $t$, always accompanied by one of these
integrable carriers.

The same variation-of-constants induction, with these carriers and
finite polynomial factors in $t$, bounds all lower-jet products through
the fixed order $N$. At the moving endpoint, $H\partial_HL=-1$;
the boundary terms retain
$e^{-L/2}$, $e^{-L}$, or $(\varphi_0+\varphi_0^2)w(L)$, so the same
induction includes them.

For ordinary terminal-height derivatives, use the equation for $R$
rather than $\log w_e$. Its explicit height dependence is
$w_e e^{-\beta_\epsilon t}e^R$, smooth also at $w_e=0$.
Each resulting source contains $e^{-\beta_\epsilon t}$ times
controlled lower variations. Thus every reserved ordinary height
derivative is bounded jointly with all auxiliary coefficient
derivatives. The original map keeps its exact front factor $w_e$.

To construct the limiting unit, set $H=0$, hence $\varphi=0$, in
the equation for $R$ and integrate on $[0,\infty)$. The convergent
integral defines
$\Psi_0(\gamma,J,w_e;\theta)$. Comparing its solution on $[0,L]$
with the positive-$H$ solution uses only the last two carriers in
\eqref{rw:secondary:integrable-carriers}; the omitted infinite tail
is $O(e^{-L/2})$. Variation of constants proves an $O(H^{1/2})$
difference at order zero. The differentiated comparisons add finitely
many powers of $|\log H|$. Choose one $q\in(0,1/2)$ to absorb all
of them after $N$ has been fixed. This proves the stated mixed
$O(H^q)$ limits.

For $x\in\{\bar a,\nu\}$ put $Y_x=\partial_xR$. Since
$\beta_\epsilon$ is independent of $x$, the exact coefficient
variations satisfy
\[
 (Y_x)_t+wQ_wY_x=\epsilon F_x,\qquad Y_x(0)=0,\qquad
 F_{\bar a}=\frac{G\varphi w}{W^2},\quad
 F_\nu=\frac{\varphi w(1-e)(W+G)}{W^2}.
\]
Dividing this equation by $\epsilon$ gives
\[
 (Z_x)_t+wQ_wZ_x=F_x,\qquad Z_x(0)=0.
\]
This equation defines $Z_x$ also at $\epsilon=0$, and uniqueness gives
$Y_x=\epsilon Z_x$. Both forcing terms have the carrier $\varphi w$.
Their auxiliary derivatives are bounded by
$C(1+t)^{M_N}\varphi e^{-t/2}$; ordinary height derivatives replace
the factor $w_e$ by bounded variations. The other differentiated
sources contain a lower variation of $Z_x$ multiplied by an integrable
variation of $wQ_w$. The same variation-of-constants induction through
total order $N-1$, including ordinary $\epsilon$ derivatives and the
moving endpoint, therefore gives
\[
 S_x:=Z_x(L),\qquad
 \|S_x\|_{N-1}^{\rm aux}
 \le C H^{1/2}(1+L)^{M_N}\le C H^q.
\]
Thus the coefficient variations have the stated factorization with
the full auxiliary derivative budget.

The terminal-height limit is nonconstant. At $\gamma=J=0$ the
limiting equation has $Q=(k+2w)/(b+w)$; direct separation gives
\[
 t=\frac bk\log\frac{w_e}{w}
       +\frac1{2k}\log\frac{k+2w}{k+2w_e},\qquad
 \Psi_0(0,0,w_e;\theta)
       =-\frac1{2b}\log(1+2w_e/k).
\]
This also verifies its ordinary smoothness at zero height.
No assertion here concerns an ordinary $C^2$ extension in $H$ at zero.

\subsection{Arrival and confinement for returning upper arcs}
\label{rw:app:secondary-barriers}

To connect the original axis input to the long map, use
$t=1/X$, $v=Z/X$ and divide the field by $X>0$.
The exact equations are
\[
\begin{aligned}
 t'/t={}&-\theta-\bar m v^2+
 t[(1+2\epsilon\bar m)v-\epsilon\bar a]
                  -\epsilon(1+\epsilon\bar m)t^2,\\
 v'={}&bv+\bar d v^2-\bar m v^3
 +t[(1+2\epsilon\bar m)v^2
                -\epsilon(\bar d+\nu+\bar a)v]\\
 &+t^2[-\epsilon(1+\epsilon\bar m)v+\epsilon^2\nu].
\end{aligned}
\]
The upper half-plane has $v\ge\epsilon t$. Fix
$V=1/\varphi_0$. For $0<v\le V$ and sufficiently small entering
$t,\epsilon,\bar d,\bar m,\nu$, these equations give $v'\ge v/2$.
The apparent constant forcing, divided by $v$, is bounded by
$\epsilon|\nu|t$. In the first equation discard the favorable
$-\theta$ and use $\epsilon|\bar a|t/v\le|\bar a|$.
Then
\[
 \frac{\dd\log t}{\dd v}\le C(1+|\bar m|v+t).
\]
Integrate on a small fixed $t$ strip and use first exit: an entry
$t=1/p$ has $t\le C_V/p$ up to $v=V$. It cannot exit through
$Y=0$, since $Y_\xi=\epsilon X>0$ there, and cannot stay forever
below $V$ because $v'\ge v/2$. Thus it reaches the fixed section
with $w=vt\to0$ as $p\to\infty$. The flight time is finite for
each entry and need not be uniformly bounded. Reflection gives
the incoming assertion.

For the middle passage choose $w_1=w_0+12$ and then $L\ge10$,
large enough for the reflected adjoining signs on $0\le w\le w_1$.
Only after these choices reduce the coefficient and scale bounds.
Put
\[
 R_0=L^2w_0,\quad K_0=L+2,\quad A_0=3(L+2),\quad
 R_*=4(L^2+L+2),\quad
 R_{\max}=R_0+2A_0L+R_*+2 .
\]
Then $R_{\max}/L^2=w_0+10+16/L+10/L^2<w_1$.
On $-L\le\hat\varphi\le L$, $0\le R\le R_{\max}$ the exact
projective field \eqref{rw:secondary:projective} converges to
\[
 \hat\varphi'=M-\hat\varphi^2-C\hat\varphi-R,\qquad
 R'=-R(\hat\varphi+C).
\]
For the normalized faces $|C|,|M|\le1$, a common scale restriction
therefore gives
\[
 R'\le K_0R,\qquad \hat\varphi'\le-R/2\quad(R\ge R_*).
\]
For $\mathcal V=R+A_0(\hat\varphi+L)$ it follows that
$\mathcal V'\le-K_0R/2<0$ when $R\ge R_*$.
At a first crossing of
$\mathcal V=R_0+2A_0L+R_*+1$, one would have $R>R_*$ and a
negative crossing derivative, a contradiction. This barrier lies
strictly below the previously fixed radial boundary.
The field points inward at $\hat\varphi=L$, and $R=0$ is invariant.
Hence an arc stays in this compact rectangle until its first meeting
with $\hat\varphi=-L$.

This is a spatial statement. A given returning upper arc must make
that meeting, because its later incoming endpoint has
$\hat\varphi=-q/(\epsilon\lambda)<-L$. It is not asserted that
every initial point in the rectangle returns. At the meeting,
$w=R/L^2<w_1$. The reflected adjoining signs then continue the
same arc outward to the fixed horizontal section with decreasing
$w$; uniqueness excludes a second middle excursion. No upper bound
for the intervening time is used.

\subsection{Relative repeated-root limits at reciprocal input zero}
\label{rw:app:secondary-tails}

The repeated upper action may diverge, while the input derivatives of
its logarithmic comparison retain ordinary limits in the actual
reciprocal tag $t=1/P$. We separate this constant action from the
relative units. In the horizontal matching $H=\lambda L$ is fixed
and the terminal height varies, so the ordinary height estimates above
apply at reciprocal zero.

Apply Corollary~\ref{rw:transitions:relative-continuity} to
\eqref{rw:secondary:repeated}, with $q=1$, $f(v)=\mu_u+v^2$ and
the normal margins of Section~\ref{rw:upper-degenerate-preparation}.
Fix all needed orders before choosing the strip so that
$C_N\varpi<\beta_*/8$. At its fixed positive exit $v_R$ the scalar
velocity has a positive margin. The full strip between $\pm v_R$
passes only for $\mu_u>0$; its time tends to infinity as the gap
vanishes. Its time derivatives have polynomial bounds in $1+T$,
and moving entries satisfy the same fixed-exit jet estimate, by
Appendix~\ref{rw:app:relative-variations}. The prepared field,
sections and reserved flow jets converge on bounded time intervals.
The corollary therefore gives convergence of the relative amplitudes
and their ordinary radial derivatives through normal input zero and
a lost exit. The absolute action remains in its constant multiplier.
For $C=2g$, $M=-g^2$, put
\[
 V=\hat z+g^{-2},\qquad Z_1=\hat x+g\hat z-g^{-1},\qquad r=Z_1/V.
\]
The limiting affine equations give
$V'=gV+Z_1$, $Z_1'=gZ_1$, $r'=-r^2$, and the first integral
$\log(Vr)-g/r$ on the relevant positive branch.
At $(P,0)$ it is
$\log(P-P_c)-P_c/(P-P_c)$, $P_c=1/g$.
The incoming and outgoing relative logarithmic derivatives are
$-P/(P-P_c)^2$ and $-1/P$. These are the limits used in
\eqref{rw:secondary:repeated-limit}.
To extend the derivative limit to reciprocal zero we need an ordinary
input bound. In the two-infinite repeated row, $\mu_u>0$ and the slow
sections are fixed.
Reserve four ordinary input orders in the tag and fixed connecting maps.
Lemma~\ref{rw:secondary:tags} gives a regular change between the actual
radial section coordinate and $t=1/P$, through zero. The connections
to the fixed slow sections stay away from the angular root. Their
normal entry is therefore $t$ times a positive unit with three bounded
ordinary derivatives. Proposition~\ref{rw:relative-core}, including
normal input zero, gives the same relative bounds across the long slow
strip. At the horizontal connections use Proposition~\ref{rw:exact-long}
with $H=\lambda L$ fixed while the terminal height varies. Its ordinary
height jets through zero supply the remaining relative units. Thus
the inputs in this row are
\[
 u_+=c_+tV_+(t),\qquad u_-=c_-tV_-(t),\qquad
 c_+=\Lambda,\quad c_-=\Lambda e^{-A_u},\quad
 0<c_-\le c_+\le1,
\]
where $V_\pm$ have common ordinary $C^3$ bounds and positive lower
margins. These bounds remain uniform as $\mu_u\downarrow0$.
Put $l_\pm=1+tV_\pm'/V_\pm\ge1/2$ on a common small interval.
Apply Lemma~\ref{rw:joint-jet-budget} with $N=2$ on the fixed coefficient
fibre, using $s=tV_+(t)$, $s'=V_+l_+>c$,
$k=(V_-/V_+)\circ t(s)$ and $1+sk'/k=l_-/l_+>c$.
The two lower inputs are $c_+s$ and
$c_-sk(s)$. Absorb $c_\pm^{h_I/\epsilon}$ into their constant
multipliers; $P_\pm(c_\pm s)$ have the required ordinary $C^3$ bounds,
and $H_\pm=h_I+uP_\pm'>c$. The lemma bounds
$\mathcal K-[\epsilon\log(K_-/K_+)-h_IA_u]$ in ordinary $C^2$,
so it bounds $\mathcal K(t)-\mathcal K(0)$ in the same norm.
Only input derivatives are taken: $A_u$ and the divergent lower actions
remain in the constants. The estimate holds on
the passing side, then restricts to the lower whole-path and
positive-central interval; the varying terminal height requires no
ordinary $C^2$ extension in $H$. The common $C^2$ bound makes the
first derivatives at $t=0$ and at a fixed small $t=\eta$ differ by
at most $C\eta$. Convergence at $\eta>0$, followed by
$\eta\downarrow0$, gives the endpoint derivative limit.
